% ============================================================
%  APPENDIX A: MATHEMATICAL BACKGROUND
%  Reference appendix for readers who need foundational review.
% ============================================================

\chapter{Математические основы}
\label{app:math}

\begin{quote}
\textit{«Чудо о том, что язык математики так удивительно пригоден для формулировки законов физики,~--- это прекрасный дар, которого мы ни понять, ни заслужить не можем».}
\par\noindent\hfill---Юджин Вигнер, 1960
\end{quote}

\bigskip

Данное приложение содержит математические основы, предполагаемые на протяжении всей книги. Читатели с прочной базой бакалавриата в области математического анализа, линейной алгебры и теории вероятностей могут обращаться к нему как к справочнику при встрече с термином, используемым без определения. Читатели, впервые встречающие эти идеи, найдут здесь достаточно подробностей, чтобы следить за основным текстом, однако им рекомендуется дополнительно изучить специализированный учебник для более глубокого освоения.

% ============================================================
\section{Основы линейной алгебры}
\label{app:linalg}
\index{linear algebra}
% ============================================================

\subsection*{Векторы и матрицы}

\textbf{Вектор} $x \in \R^n$~--- это упорядоченный набор из $n$ вещественных чисел. По умолчанию мы пишем столбцовые векторы; $x^\top$ обозначает строковый транспонированный. \textbf{Скалярное произведение} $x, y \in \R^n$ равно
\[
    \langle x, y \rangle = x^\top y = \sum_{i=1}^n x_i y_i.
\]
\textbf{Матрица} $A \in \R^{m \times n}$ отображает $\R^n \to \R^m$ посредством $y = Ax$. \textbf{Транспонированная} $A^\top \in \R^{n \times m}$ удовлетворяет $(A^\top)_{ij} = A_{ji}$.

\begin{definition}[Нормы]
\label{def:norms}
$\ell^p$-норма вектора $x \in \R^n$ равна
\[
    \|x\|_p = \Bigl(\sum_{i=1}^n |x_i|^p\Bigr)^{1/p}, \qquad p \geq 1.
\]
Частные случаи: $\|x\|_1 = \sum_i |x_i|$; $\|x\|_2 = \sqrt{\sum_i x_i^2}$ (евклидова норма, записываемая $\|x\|$ в тех случаях, когда нет неоднозначности); $\|x\|_\infty = \max_i |x_i|$. \textbf{Норма Фробениуса} матрицы равна $\|A\|_F = \sqrt{\sum_{i,j} A_{ij}^2}$.
\end{definition}

\subsection*{Собственные значения и спектральные свойства}

Для квадратной матрицы $A \in \R^{n \times n}$ скаляр $\lambda$ и ненулевой вектор $v$, удовлетворяющие $Av = \lambda v$, называются \textbf{собственным значением}\index{eigenvalue} и \textbf{собственным вектором}. Совокупность всех собственных значений называется \textbf{спектром} матрицы $A$.

\begin{definition}[Спектральный радиус]
\label{def:spectral-radius}
Спектральный радиус матрицы $A$ равен $\rho(A) = \max_i |\lambda_i|$, где $\lambda_1, \ldots, \lambda_n$~--- собственные значения $A$.
\end{definition}

Матрица $A$ называется \textbf{положительно полуопределённой} (ПП) если $x^\top A x \geq 0$ для всех $x$, что равносильно тому, что все собственные значения неотрицательны. ПП-матрицы естественно возникают как ковариационные матрицы и гессианы в точках локального минимума.

\begin{example}[Ковариационная матрица]
Пусть $X \in \R^{n \times d}$~--- матрица данных со столбцами с нулевым средним. Эмпирическая ковариация $\Sigma = X^\top X / n$ является ПП. Её собственные векторы~--- это главные компоненты, используемые в PCA, а собственные значения измеряют дисперсию вдоль каждой компоненты.
\end{example}

\subsection*{Линейные системы и обратимость}

Квадратная матрица $A$ является \textbf{обратимой} (невырожденной) тогда и только тогда, когда $\det(A) \neq 0$, что равносильно тому, что все её собственные значения ненулевые. Если $A$ обратима, единственное решение системы $Ax = b$ равно $x = A^{-1}b$.

В обучении с подкреплением системы вида $(I - \gamma P^\pi) V = R$ возникают из уравнений Беллмана (глава~\ref{ch:dp}). Поскольку $\rho(\gamma P^\pi) = \gamma < 1$, матрица $I - \gamma P^\pi$ обратима, что гарантирует единственность функции ценности.

% ============================================================
\section{Математический анализ и оптимизация}
\label{app:calculus}
\index{calculus}\index{optimisation}
% ============================================================

\subsection*{Градиенты и гессианы}

Пусть $f : \R^n \to \R$ дифференцируема. \textbf{Градиент}\index{gradient}~--- это вектор частных производных:
\[
    \nabla_\theta f(\theta)
    = \Bigl(\frac{\partial f}{\partial \theta_1}, \ldots,
             \frac{\partial f}{\partial \theta_n}\Bigr)^\top \in \R^n.
\]
Градиент указывает в направлении наибольшего возрастания. \textbf{Гессиан}\index{Hessian}~--- это матрица вторых производных:
\[
    H(\theta) = \nabla^2 f(\theta), \qquad
    H_{ij} = \frac{\partial^2 f}{\partial \theta_i \partial \theta_j}.
\]
Если $f$ дважды непрерывно дифференцируема, то $H$ симметрична.

\begin{definition}[Критическая точка]
$\theta^*$ является \textbf{критической точкой} $f$, если $\nabla f(\theta^*)=0$. Это точка локального минимума, если $H(\theta^*)$ является ПП, и точка локального максимума, если $-H(\theta^*)$ является ПП.
\end{definition}

\subsection*{Правило цепочки}

Если $g: \R^m \to \R^n$ и $f: \R^n \to \R$, то правило цепочки даёт
\[
    \nabla_\theta (f \circ g)(\theta)
    = J_g(\theta)^\top \nabla_{g(\theta)} f,
\]
где $J_g \in \R^{n \times m}$~--- якобиан функции $g$. Автоматическое дифференцирование в системах глубокого обучения эффективно вычисляет это выражение посредством обратного распространения ошибки.

\subsection*{Градиентный спуск}

\textbf{Градиентный спуск}\index{gradient descent} итеративно выполняет
\[
    \theta_{t+1} = \theta_t - \alpha \nabla f(\theta_t),
\]
где $\alpha > 0$~--- \textbf{скорость обучения}. Обновление перемещает $\theta$ в направлении наибольшего убывания.

\textbf{Стохастический градиентный спуск} (SGD) заменяет истинный градиент зашумлённой оценкой $\hat{g}_t$, вычисленной на минибатче:
\[
    \theta_{t+1} = \theta_t - \alpha \hat{g}_t, \qquad
    \E[\hat{g}_t] = \nabla f(\theta_t).
\]
Дисперсия $\hat{g}_t$ выступает как неявная регуляризация и может помогать выходить из острых локальных минимумов.

\textbf{Adam}\index{Adam optimizer} (метод адаптивных моментов) поддерживает экспоненциальные скользящие средние градиента $m_t$ и квадрата градиента $v_t$:
\begin{align*}
    m_t &= \beta_1 m_{t-1} + (1-\beta_1)\hat{g}_t, \\
    v_t &= \beta_2 v_{t-1} + (1-\beta_2)\hat{g}_t^2, \\
    \theta_{t+1} &= \theta_t - \frac{\alpha}{\sqrt{\hat{v}_t}+\epsilon}\,\hat{m}_t,
\end{align*}
где $\hat{m}_t = m_t/(1-\beta_1^t)$ и $\hat{v}_t = v_t/(1-\beta_2^t)$~--- скорректированные оценки. Значения по умолчанию $\beta_1=0.9$, $\beta_2=0.999$, $\epsilon=10^{-8}$ хорошо работают в широком круге задач.

\subsection*{Выпуклость и сходимость}

\begin{definition}[Выпуклость]
\label{def:convexity}
Функция $f: \R^n \to \R$ является \textbf{выпуклой}, если для всех $x,y$ и $\lambda \in [0,1]$:
\[
    f(\lambda x + (1-\lambda)y) \leq \lambda f(x) + (1-\lambda)f(y).
\]
Она является \textbf{сильно выпуклой} с параметром $\mu > 0$, если $f(x) - \frac{\mu}{2}\|x\|^2$ выпукла.
\end{definition}

Для выпуклой функции с $L$-липшицевым градиентом (т.е. $\|\nabla f(x)-\nabla f(y)\| \leq L\|x-y\|$) градиентный спуск с размером шага $\alpha \leq 1/L$ сходится к глобальному минимуму. Скорость сходимости составляет $O(1/t)$ для выпуклых функций и $O(e^{-\mu t/L})$ для сильно выпуклых.

\begin{example}[Минимизация квадратичной функции]
Пусть $f(\theta) = \frac{1}{2}\theta^\top A\theta - b^\top\theta$ при ПП матрице $A$. Тогда $\nabla f(\theta) = A\theta - b$, и обновление градиентного спуска принимает вид $\theta_{t+1} = \theta_t - \alpha(A\theta_t - b)$. Единственный минимум равен $\theta^* = A^{-1}b$. При $\alpha < 2/\lambda_{\max}(A)$ итераты сходятся геометрически.
\end{example}

% ============================================================
\section{Теория вероятностей}
\label{app:probability}
\index{probability theory}
% ============================================================

\subsection*{Случайные величины и распределения}

\textbf{Случайная величина} $X$~--- это измеримая функция из вероятностного пространства $\Omega$ в $\R$. Для дискретной $X$ \textbf{функция масс вероятности} (ФМВ) равна $p(x) = \Prob(X=x)$. Для непрерывной $X$ \textbf{функция плотности вероятности} (ФПВ) $p$ удовлетворяет $\Prob(a \leq X \leq b) = \int_a^b p(x)\,dx$.

\begin{definition}[Математическое ожидание и дисперсия]
\label{def:expectation}
\[
    \E[X] = \sum_x x\,p(x) \quad\mbox{(дискретный случай)}, \qquad
    \E[X] = \int x\,p(x)\,dx \quad\mbox{(непрерывный случай)}.
\]
\[
    \Var[X] = \E\!\left[(X - \E[X])^2\right] = \E[X^2] - (\E[X])^2.
\]
\textbf{Ковариация} $X$ и $Y$ равна $\Cov[X,Y] = \E[(X-\E X)(Y-\E Y)]$.
\end{definition}

Математическое ожидание линейно: $\E[aX + bY] = a\E[X] + b\E[Y]$ для любых констант $a, b$ и случайных величин $X, Y$ (независимо от их зависимости). Дисперсия квадратична: $\Var[aX] = a^2 \Var[X]$, и для независимых $X, Y$ выполняется $\Var[X+Y] = \Var[X] + \Var[Y]$.

\subsection*{Условная вероятность и теорема Байеса}

\textbf{Условная вероятность} события $A$ при условии $B$ равна
\[
    \Prob(A \mid B) = \frac{\Prob(A \cap B)}{\Prob(B)}, \quad \Prob(B) > 0.
\]

\begin{theorem}[Теорема Байеса]
\label{thm:bayes}
\[
    \Prob(A \mid B) = \frac{\Prob(B \mid A)\,\Prob(A)}{\Prob(B)}.
\]
В непрерывном случае с априорным распределением $p(\theta)$ и правдоподобием $p(x|\theta)$:
\[
    p(\theta \mid x) = \frac{p(x \mid \theta)\,p(\theta)}{p(x)}.
\]
\end{theorem}

Теорема Байеса лежит в основе байесовских стратегий исследования (апостериорное сэмплирование, глава~\ref{ch:bandits}) и моделирования вознаграждения в RLHF (глава~\ref{ch:practical-rlhf}).

\subsection*{Условное математическое ожидание и свойство башни}

\textbf{Условное математическое ожидание} $\E[X \mid Y]$~--- само по себе случайная величина\,---\,функция от $Y$. Его значение при $Y=y$ равно $\E[X \mid Y=y] = \int x\,p(x \mid y)\,dx$.

\begin{theorem}[Свойство башни / Закон полного математического ожидания]
\label{thm:tower}
\[
    \E[X] = \E\!\bigl[\E[X \mid Y]\bigr].
\]
Более общо, для $\sigma$-алгебр $\mathcal{F}_1 \subseteq \mathcal{F}_2$: $\E[\E[X|\mathcal{F}_2]|\mathcal{F}_1] = \E[X|\mathcal{F}_1]$.
\end{theorem}

Свойство башни используется в RL всякий раз, когда мы маргинализируем по будущим траекториям или вычисляем ожидаемый доход путём обусловливания на текущем состоянии.

\subsection*{Концентрационные неравенства}

\begin{theorem}[Закон больших чисел]
\label{thm:lln}
Пусть $X_1, X_2, \ldots$~--- н.о.р. с математическим ожиданием $\mu$. Тогда при $n \to \infty$:
\[
    \bar{X}_n = \frac{1}{n}\sum_{i=1}^n X_i \xrightarrow{\,\mbox{п.н.}\,} \mu.
\]
\end{theorem}

ЗБЧ обосновывает использование выборочных средних для оценки ожидаемого дохода в методах Монте-Карло (глава~\ref{ch:mc}).

\begin{theorem}[Центральная предельная теорема]
\label{thm:clt}
При тех же условиях, что и в Теореме~\ref{thm:lln}, и при $\sigma^2 = \Var[X_i] < \infty$:
\[
    \sqrt{n}\,(\bar{X}_n - \mu) \xrightarrow{\,d\,} \cN(0, \sigma^2).
\]
\end{theorem}

ЦПТ даёт доверительные интервалы для оценок Монте-Карло и мотивирует гауссовскую модель шума в непрерывном управлении.

\begin{example}[Оценка среднего дохода]
Предположим, что мы проводим $n$ независимых эпизодов и записываем доходы $G^{(1)}, \ldots, G^{(n)}$. По ЗБЧ $\bar{G}_n \to V^\pi(s_0)$. По ЦПТ доверительный интервал уровня 95\% приблизительно равен $\bar{G}_n \pm 1.96\,\hat{\sigma}/\sqrt{n}$, где $\hat{\sigma}^2$~--- выборочная дисперсия доходов.
\end{example}

% ============================================================
\section{Методы выборки}
\label{app:sampling}
\index{sampling methods}
% ============================================================

Многие величины в RL\,---\,ожидаемые доходы, градиенты стратегии, KL-штрафы\,---\,не поддаются аналитическому вычислению, но могут быть \emph{оценены} по выборкам. В этом разделе рассматриваются основные стратегии выборки.

\subsection*{Оценка методом Монте-Карло}

Если математическое ожидание $\mu = \E_{X \sim p}[f(X)]$ трудно вычислить аналитически, \textbf{оценка Монте-Карло}\index{Monte Carlo} определяется как
\[
    \hat{\mu}_n = \frac{1}{n}\sum_{i=1}^n f(x^{(i)}), \qquad
    x^{(i)} \sim p.
\]
По ЗБЧ $\hat{\mu}_n \to \mu$. Дисперсия оценки равна $\Var[\hat{\mu}_n] = \Var[f(X)] / n$, поэтому удвоение числа выборок вдвое уменьшает стандартную ошибку.

Методы Монте-Карло в RL (глава~\ref{ch:mc}) используют полные доходы эпизодов как несмещённые (хотя и высокодисперсные) оценки функции ценности.

\subsection*{Выборка по значимости}

Когда выборки из $p$ недоступны, но доступны выборки из другого \textbf{предложения} $q$, \textbf{выборка по значимости}\index{importance sampling} корректирует это с помощью отношения правдоподобий:
\[
    \mu = \E_{X \sim p}[f(X)]
        = \E_{X \sim q}\!\left[\frac{p(X)}{q(X)}\,f(X)\right].
\]
Отношение $w(x) = p(x)/q(x)$ называется \textbf{весом значимости}. Оценка несмещена при условии $q(x) > 0$ там, где $p(x)f(x) \neq 0$, однако может иметь высокую (и даже бесконечную) дисперсию, если $p$ и $q$ сильно различаются.

Выборка по значимости применяется в обучении TD по внешней стратегии и занимает центральное место в целевой функции PPO (глава~\ref{ch:actor-critic}), где $q = \pi_{\theta_{\text{old}}}$ и $p = \pi_\theta$. Усечение отношения $\rho_t = \pi_\theta / \pi_{\theta_{\text{old}}}$ контролирует дисперсию.

\subsection*{Выборка с отклонением}

\textbf{Выборка с отклонением}\index{rejection sampling} позволяет извлекать выборки из целевого распределения $p(x) \propto \tilde{p}(x)$ с помощью предложения $q(x)$ такого, что $\tilde{p}(x) \leq M q(x)$ для всех $x$:

\begin{enumerate}[leftmargin=*]
    \item Извлечь $x \sim q$.
    \item Принять $x$ с вероятностью $\tilde{p}(x)/(M q(x))$; в противном случае отклонить.
\end{enumerate}

Вероятность принятия равна $1/M$, поэтому «плотное» предложение (малое $M$) критически важно для эффективности. Выборка с отклонением используется при анализе офлайновых моделей предпочтений (глава~\ref{ch:practical-rlhf}) и в некоторых схемах структурированного декодирования.

% ============================================================
\section{Теория информации}
\label{app:info-theory}
\index{information theory}
% ============================================================

Теория информации количественно описывает неопределённость и стоимость передачи данных. В RL её концепции появляются при регуляризации вознаграждения, в задачах оптимизации стратегии и в анализе исследования.

\begin{definition}[Энтропия]
\label{def:entropy}
\textbf{Энтропия Шеннона}\index{entropy} дискретного распределения $p$ над алфавитом $\cV$ равна
\[
    H(p) = -\sum_{x \in \cV} p(x)\log p(x),
\]
измеряемая в натах (основание $e$) или битах (основание 2). Энтропия максимальна для равномерного распределения и равна нулю для точечной меры.
\end{definition}

Бонусы за энтропию $\mathcal{H}(\pi(\cdot|s))$ стимулируют исследование, штрафуя стратегии, вырождающиеся в детерминированные действия (глава~\ref{ch:actor-critic}).

\begin{definition}[Перекрёстная энтропия]
\label{def:cross-entropy}
\textbf{Перекрёстная энтропия} между распределениями $p$ и $q$ равна
\[
    H(p,q) = -\sum_x p(x)\log q(x).
\]
Она равна $H(p) + \KL{p}{q}$ и является стандартной целевой функцией при предобучении языковых моделей и дообучении с учителем.
\end{definition}

\begin{definition}[KL-дивергенция]
\label{def:kl}
\textbf{Дивергенция Кульбака-Лейблера}\index{KL divergence} от $q$ к $p$ равна
\[
    \KL{p}{q} = \sum_x p(x)\log\frac{p(x)}{q(x)} \geq 0,
\]
причём равенство выполняется тогда и только тогда, когда $p = q$ почти всюду. Заметим, что $\KL{p}{q} \neq \KL{q}{p}$ в общем случае; KL-дивергенция не является метрикой.
\end{definition}

KL-штраф $\beta\,\KL{\pi_\theta}{\pi_\mathrm{ref}}$ в задачах RLHF (глава~\ref{ch:practical-rlhf}) и DPO (глава~\ref{ch:lm-post-training}) штрафует за чрезмерное отклонение от эталонной стратегии, предотвращая взлом вознаграждения.

\begin{definition}[Взаимная информация]
\label{def:mutual-info}
\textbf{Взаимная информация} между случайными величинами $X$ и $Y$ равна
\[
    I(X;Y) = \KL{p(x,y)}{p(x)p(y)}
            = H(X) - H(X \mid Y) \geq 0.
\]
Она измеряет уменьшение неопределённости об $X$ после наблюдения $Y$.
\end{definition}

\begin{example}[RL с KL-регуляризацией]
Оптимальная стратегия для KL-регуляризованной задачи $\max_\pi \E_\pi[R] - \beta\,\KL{\pi}{\pi_\mathrm{ref}}$ имеет замкнутую форму: $\pi^*(a|s) \propto \pi_\mathrm{ref}(a|s)\, e^{Q^*(s,a)/\beta}$, где $Q^*$~--- мягкая Q-функция. Это связано с выводом DPO в главе~\ref{ch:dpo}.
\end{example}

% ============================================================
\section{Операторы и сжимающие отображения}
\label{app:operators}
\index{contraction mapping}\index{Bellman operator}
% ============================================================

Доказательства сходимости динамического программирования и TD-обучения опираются на классический результат функционального анализа~--- теорему Банаха о неподвижной точке. В этом разделе развивается необходимая база и показывается, как оператор Беллмана вписывается в данную схему.

\subsection*{Метрические пространства и сжимающие отображения}

\begin{definition}[Метрическое пространство]
\label{def:metric-space}
\textbf{Метрическое пространство} $(X, d)$~--- это множество $X$, снабжённое функцией расстояния $d: X \times X \to [0,\infty)$, удовлетворяющей: (i) $d(x,y)=0 \Leftrightarrow x=y$; (ii) $d(x,y)=d(y,x)$; (iii) $d(x,z) \leq d(x,y)+d(y,z)$ (неравенство треугольника). Метрическое пространство называется \textbf{полным}, если каждая последовательность Коши сходится в $X$.
\end{definition}

\begin{definition}[Сжимающее отображение]
\label{def:contraction}
Оператор $T: X \to X$ на метрическом пространстве $(X,d)$ является \textbf{сжимающим}\index{contraction} с модулем $\kappa \in [0,1)$, если
\[
    d(Tx, Ty) \leq \kappa\, d(x, y) \quad \mbox{для всех } x, y \in X.
\]
\end{definition}

\begin{theorem}[Теорема Банаха о неподвижной точке]
\label{thm:banach}
Если $T$~--- сжимающее отображение на полном метрическом пространстве $(X,d)$, то $T$ имеет \emph{единственную} неподвижную точку $x^*$ (т.е. $Tx^* = x^*$), и для любой начальной точки $x_0 \in X$ итераты $x_{k+1} = T x_k$ сходятся:
\[
    d(x_k, x^*) \leq \frac{\kappa^k}{1-\kappa}\,d(x_1, x_0).
\]
Сходимость геометрическая со скоростью $\kappa$.
\end{theorem}

\subsection*{Оператор Беллмана как сжимающее отображение}

Обозначим через $\mathcal{B}(\cS)$ пространство ограниченных вещественнозначных функций на $\cS$, снабжённое sup-нормой $\|V\|_\infty = \max_s |V(s)|$. Это полное метрическое пространство.

\begin{definition}[Операторы Беллмана]
\label{def:bellman-op}
Для фиксированной стратегии $\pi$ \textbf{оператор оценки Беллмана} $\mathcal{T}^\pi$ отображает $V : \cS \to \R$ в
\[
    (\mathcal{T}^\pi V)(s)
    = \sum_a \pi(a|s)\Bigl[R(s,a) + \gamma \sum_{s'} P(s'|s,a)\,V(s')\Bigr].
\]
\textbf{Оператор оптимальности Беллмана} $\mathcal{T}^*$ отображает $V$ в
\[
    (\mathcal{T}^* V)(s)
    = \max_a \Bigl[R(s,a) + \gamma \sum_{s'} P(s'|s,a)\,V(s')\Bigr].
\]
\end{definition}

\begin{theorem}[Сжимающее свойство Беллмана]
\label{thm:bellman-contraction}
Оба оператора $\mathcal{T}^\pi$ и $\mathcal{T}^*$ являются сжимающими на $(\mathcal{B}(\cS), \|\cdot\|_\infty)$ с модулем $\gamma$.
\end{theorem}

\begin{proof}
Докажем результат для $\mathcal{T}^*$; доказательство для $\mathcal{T}^\pi$ аналогично. Для любых $V, W \in \mathcal{B}(\cS)$ и любого состояния $s$:
\begin{align*}
|(\mathcal{T}^* V)(s) - (\mathcal{T}^* W)(s)|
&= \Bigl|\max_a [R(s,a) + \gamma(PV)(s,a)]
   - \max_a [R(s,a) + \gamma(PW)(s,a)]\Bigr| \\
&\leq \max_a \gamma\,\Bigl|\sum_{s'} P(s'|s,a)[V(s')-W(s')]\Bigr| \\
&\leq \gamma\,\|V - W\|_\infty,
\end{align*}
где мы воспользовались неравенством $|\max_a f(a) - \max_a g(a)| \leq \max_a |f(a)-g(a)|$ и тем, что $P(\cdot|s,a)$~--- распределение вероятностей. Взяв супремум по $s$, получаем $\|\mathcal{T}^* V - \mathcal{T}^* W\|_\infty \leq \gamma\,\|V-W\|_\infty$. Так как $\gamma < 1$, $\mathcal{T}^*$ является сжимающим.
\end{proof}

\begin{corollary}[Существование и единственность функций ценности]
\label{cor:value-exists}
Для любого МДП с $\gamma < 1$ (или конечным горизонтом) функции ценности $V^\pi$ и $V^*$ существуют и являются единственными неподвижными точками операторов $\mathcal{T}^\pi$ и $\mathcal{T}^*$ соответственно. Начиная с любой ограниченной $V_0$, итераты $V_{k+1} = \mathcal{T}^\pi V_k$ (оценка стратегии) и $V_{k+1} = \mathcal{T}^* V_k$ (итерация ценности) сходятся геометрически:
\[
    \|V_k - V^\pi\|_\infty \leq \gamma^k\,\|V_0 - V^\pi\|_\infty.
\]
\end{corollary}

Этот результат обосновывает алгоритмы, разработанные в главах~\ref{ch:dp} и~\ref{ch:td}: каждый шаг оценки стратегии или итерации ценности уменьшает ошибку аппроксимации в $\gamma$ раз.

\begin{example}[Итерация ценности на двухсостоянии МДП]
Рассмотрим МДП с двумя состояниями $\cS = \{s_1, s_2\}$, $\gamma = 0.9$ и оптимальными значениями $V^* = (10, 5)^\top$. Начиная с $V_0 = (0,0)^\top$, после $k$ итераций ошибка удовлетворяет $\|V_k - V^*\|_\infty \leq 0.9^k \cdot 10$. Для достижения ошибки $\leq 0.01$ необходимо $k \geq \log(0.001)/\log(0.9) \approx 66$ итераций~--- точная гарантия сходимости, не зависящая от структуры вознаграждения.
\end{example}

\subsection*{Нерасширяющие операторы и приближённое ДП}

При использовании аппроксимации функций (глава~\ref{ch:value-approx}) проецированный оператор Беллмана $\Pi \mathcal{T}^\pi$ (где $\Pi$~--- ортогональная проекция на пространство аппроксимаций) в общем случае больше не является сжимающим. Теорема о неподвижной точке TD даёт более слабую, но всё же полезную гарантию: при линейной аппроксимации функций и определённых условиях выборки по собственной стратегии TD сходится к неподвижной точке, ошибка в которой ограничена кратным наилучшей достижимой ошибки аппроксимации.

% ============================================================
%  End of Appendix A
% ============================================================
